\subsection{The small-prime hypothesis in the exponent-thirteen argument}
\label{subsec:thirteen-small-prime-gate}

Let $A_B$ be the genus-two Jacobian in the exponent-thirteen
decomposition, let $\chi$ be a cubic Dirichlet character modulo
thirteen, and put $\mathcal O=\mathbb Z_3[\zeta_3]$ with maximal
ideal $\mathfrak m=(1-\zeta_3)$.  The large split-prime argument
does not apply at three merely because three splits in the quartic
CM field.

\begin{proposition}[A small-prime hypothesis obstruction]
\label{prop:thirteen-small-prime-mr}
For $a=1,2$, the natural lattice
\[
 T=(T_3A_B\otimes_{\mathbb Z_3}\mathcal O)\otimes\chi^a
\]
does not satisfy hypothesis (H.4) of
\citet[\S3.5, p.~27]{MazurRubin2004}.  The same holds for its
dual and for every integral Tate twist of either lattice.
This is a limitation on that theorem's hypotheses, not a
nonvanishing statement for any Selmer group.
\end{proposition}
\begin{proof}
The residue field of $\mathcal O$ is $\mathbb F_3$, and the
cubic character becomes trivial modulo $\mathfrak m$.  Consequently
$T/\mathfrak mT\simeq A_B[3]$ as $\mathbb F_3[G_{\mathbb Q}]$-modules.
The principal polarization of the Jacobian gives a perfect
Galois-equivariant alternating Weil pairing
\[
 A_B[3]\times A_B[3]\longrightarrow\mu_3.
\]
It identifies this residual module with its Cartier dual and
therefore gives a nonzero element of
\[
 \operatorname{Hom}_{\mathbb F_3[G_{\mathbb Q}]}
       (T/\mathfrak mT,T^*[\mathfrak m]).
\]
Thus (H.4a) fails.  The alternative (H.4b) requires $p>4$, so
it fails at $p=3$ as well.

The mod-three cyclotomic character has order two.  Multiplying
a residual representation by any of its powers preserves
Cartier self-duality, since the square of the twisting character
is one.  Moreover the ordinary dual of $A_B[3]$ is its
cyclotomic twist.  This proves the assertion for duals and
integral Tate twists.
\end{proof}

The proposition concerns the displayed natural lattices.  It does
not assert that every stable lattice in the rational representation
has the same extension structure.  It does show why the verified
large-prime proof cannot simply be specialized to three.  A
different small-prime argument, or a theorem with different
hypotheses, would be needed to obtain the missing three-primary
Selmer input by this route.
